% ============================================================================
% CAPÍTULO 11: INTEGRATED ANALYSIS
% Bioinformática para Biologia de Sistemas
% ============================================================================

% ============================================================================
% TEMPLATE BASE PARA APRESENTAÇÕES EM BEAMER
% Bioinformática para Biologia de Sistemas
% ============================================================================

\documentclass[12pt, aspectratio=169]{beamer}

% ============================================================================
% PACOTES ESSENCIAIS
% ============================================================================
\usepackage[utf8]{inputenc}
\usepackage[T1]{fontenc}
\usepackage[brazilian]{babel}
\usepackage{amsmath, amssymb, amsthm}
\usepackage{graphicx}
\usepackage{xcolor}
\usepackage{tikz}
\usetikzlibrary{arrows.meta, shapes, positioning, calc, decorations.pathreplacing, decorations.shapes, decorations.pathmorphing}
\usepackage{listings}
\usepackage{booktabs}
\usepackage{multirow}
\usepackage{hyperref}
\usepackage{chemfig}
\usepackage{chemarr}

% ============================================================================
% CONFIGURAÇÃO DO TEMA
% ============================================================================
\usetheme{Madrid}
\usecolortheme{default}

% Cores personalizadas
\definecolor{azulescuro}{RGB}{0, 51, 102}
\definecolor{azulclaro}{RGB}{51, 153, 255}
\definecolor{verdeacido}{RGB}{102, 204, 0}
\definecolor{laranjacel}{RGB}{255, 153, 51}
\definecolor{roxodna}{RGB}{153, 51, 255}

\setbeamercolor{structure}{fg=azulescuro}
\setbeamercolor{palette primary}{bg=azulescuro,fg=white}
\setbeamercolor{palette secondary}{bg=azulclaro,fg=white}
\setbeamercolor{palette tertiary}{bg=azulescuro,fg=white}
\setbeamercolor{palette quaternary}{bg=azulclaro,fg=white}
\setbeamercolor{block title}{bg=azulescuro,fg=white}
\setbeamercolor{block body}{bg=azulclaro!10,fg=black}
\setbeamercolor{block title alerted}{bg=red!80,fg=white}
\setbeamercolor{block body alerted}{bg=red!10,fg=black}
\setbeamercolor{block title example}{bg=verdeacido!80,fg=white}
\setbeamercolor{block body example}{bg=verdeacido!10,fg=black}

% ============================================================================
% CONFIGURAÇÃO DE CÓDIGO
% ============================================================================
\lstset{
    language=Python,
    basicstyle=\ttfamily\small,
    keywordstyle=\color{azulescuro}\bfseries,
    commentstyle=\color{verdeacido},
    stringstyle=\color{laranjacel},
    numbers=left,
    numberstyle=\tiny\color{gray},
    stepnumber=1,
    numbersep=5pt,
    backgroundcolor=\color{white},
    showspaces=false,
    showstringspaces=false,
    showtabs=false,
    frame=single,
    rulecolor=\color{black},
    tabsize=2,
    captionpos=b,
    breaklines=true,
    breakatwhitespace=false,
    escapeinside={\%*}{*)},
    xleftmargin=10pt,
    xrightmargin=10pt
}

% Definir estilo para R
\lstdefinestyle{Rstyle}{
    language=R,
    basicstyle=\ttfamily\small,
    keywordstyle=\color{azulescuro}\bfseries,
    commentstyle=\color{verdeacido},
    stringstyle=\color{laranjacel}
}

% Definir estilo para MATLAB
\lstdefinestyle{MATLABstyle}{
    language=Matlab,
    basicstyle=\ttfamily\small,
    keywordstyle=\color{azulescuro}\bfseries,
    commentstyle=\color{verdeacido},
    stringstyle=\color{laranjacel}
}

% ============================================================================
% COMANDOS PERSONALIZADOS
% ============================================================================

% Comando para equações destacadas
\newcommand{\destaque}[1]{%
    \begin{alertblock}{}
        \begin{center}
            #1
        \end{center}
    \end{alertblock}
}

% Comando para definições
\newcommand{\definicao}[2]{%
    \begin{block}{Definição: #1}
        #2
    \end{block}
}

% Comando para exemplos
\newcommand{\exemplo}[2]{%
    \begin{exampleblock}{Exemplo: #1}
        #2
    \end{exampleblock}
}

% Comando para observações importantes
\newcommand{\importante}[1]{%
    \begin{alertblock}{Importante!}
        #1
    \end{alertblock}
}

% Comandos matemáticos úteis
\newcommand{\R}{\mathbb{R}}
\newcommand{\N}{\mathbb{N}}
\newcommand{\Z}{\mathbb{Z}}
\newcommand{\dif}{\mathrm{d}}
\newcommand{\parcial}[2]{\frac{\partial #1}{\partial #2}}
\newcommand{\derivada}[2]{\frac{\dif #1}{\dif #2}}
\newcommand{\vect}[1]{\boldsymbol{#1}}

% Comandos biológicos
\newcommand{\gene}[1]{\textit{#1}}
\newcommand{\proteina}[1]{\textsc{#1}}
\newcommand{\especie}[1]{\textit{#1}}

% ============================================================================
% CONFIGURAÇÃO DE RODAPÉ
% ============================================================================
\setbeamertemplate{footline}{
  \leavevmode%
  \hbox{%
  \begin{beamercolorbox}[wd=.33\paperwidth,ht=2.25ex,dp=1ex,center]{author in head/foot}%
    \usebeamerfont{author in head/foot}\insertshortauthor
  \end{beamercolorbox}%
  \begin{beamercolorbox}[wd=.34\paperwidth,ht=2.25ex,dp=1ex,center]{title in head/foot}%
    \usebeamerfont{title in head/foot}\insertshorttitle
  \end{beamercolorbox}%
  \begin{beamercolorbox}[wd=.33\paperwidth,ht=2.25ex,dp=1ex,right]{date in head/foot}%
    \usebeamerfont{date in head/foot}\insertshortdate{}\hspace*{2em}
    \insertframenumber{} / \inserttotalframenumber\hspace*{2ex}
  \end{beamercolorbox}}%
  \vskip0pt%
}

% ============================================================================
% CONFIGURAÇÃO DE NAVEGAÇÃO
% ============================================================================
\setbeamertemplate{navigation symbols}{}

% ============================================================================
% CONFIGURAÇÃO DE ITEMIZE/ENUMERATE
% ============================================================================
\setbeamertemplate{itemize items}[circle]
\setbeamertemplate{enumerate items}[default]

% ============================================================================
% FIM DO TEMPLATE
% ============================================================================


% ============================================================================
% INFORMAÇÕES DO DOCUMENTO
% ============================================================================
\title[Cap. 11: Integrated Analysis]{Capítulo 11: Integrated Analysis}
\subtitle{Bioinformática para Biologia de Sistemas}
\author{Ney Lemke}
\institute{DBF-Unesp}
\date{\today}

\begin{document}

% ============================================================================
% SLIDE DE TÍTULO
% ============================================================================
\begin{frame}
\titlepage
\end{frame}

% ============================================================================
% SUMÁRIO
% ============================================================================
\begin{frame}{Sumário}
\tableofcontents
\end{frame}

% ============================================================================
% SEÇÃO 1: INTRODUÇÃO
% ============================================================================
\section{Introdução}

\begin{frame}{Visão Geral do Capítulo}
\begin{block}{Objetivos de Aprendizagem}
\begin{itemize}
    \item Compreender a abordagem de biologia de sistemas
    \item Dominar estratégias de integração multi-ômica
    \item Analisar caso de estudo: ciclo da trealose em levedura
    \item Aplicar métodos de análise de vias integradas
    \item Construir modelos de sistemas integrados
\end{itemize}
\end{block}

\vspace{0.3cm}

\begin{exampleblock}{Por que Análise Integrada?}
\begin{itemize}
    \item Visão holística de sistemas biológicos complexos
    \item Descoberta de mecanismos regulatórios ocultos
    \item Predição de comportamento do sistema
    \item Identificação de alvos terapêuticos
\end{itemize}
\end{exampleblock}
\end{frame}

\begin{frame}{Abordagem de Biologia de Sistemas}
\definicao{Biologia de Sistemas}{
Estudo de sistemas biológicos através da integração quantitativa de dados multi-ômicos para compreender propriedades emergentes e comportamento do sistema como um todo.
}

\vspace{0.3cm}

\begin{center}
\begin{tikzpicture}[scale=0.8, every node/.style={font=\small}]
    % Multi-scale integration
    \node[draw, circle, fill=roxodna!30, text width=2cm, align=center] at (0,3) {Molecular};
    \node[draw, circle, fill=laranjacel!30, text width=2cm, align=center] at (3,3) {Celular};
    \node[draw, circle, fill=verdeacido!30, text width=2cm, align=center] at (6,3) {Tecido};
    \node[draw, circle, fill=azulclaro!30, text width=2cm, align=center] at (9,3) {Organismo};

    % Arrows showing integration
    \draw[->, ultra thick] (1,3) -- (2,3);
    \draw[->, ultra thick] (4,3) -- (5,3);
    \draw[->, ultra thick] (7,3) -- (8,3);

    % Bottom layer: omics
    \node[font=\tiny] at (0,1.5) {Genômica};
    \node[font=\tiny] at (2,1.5) {Transcriptômica};
    \node[font=\tiny] at (4,1.5) {Proteômica};
    \node[font=\tiny] at (6,1.5) {Metabolômica};
    \node[font=\tiny] at (8,1.5) {Fenômica};

    \draw[<->, thick, azulescuro] (0,2) -- (8,2);
    \node[above, font=\footnotesize] at (4,2) {Integração Multi-Escala};
\end{tikzpicture}
\end{center}
\end{frame}

\begin{frame}{Estratégia de Integração Multi-Ômica}
\begin{columns}
\column{0.5\textwidth}
\textbf{Camadas de Dados:}
\begin{itemize}
    \item \textbf{Genômica:} sequência, variações
    \item \textbf{Transcriptômica:} expressão gênica
    \item \textbf{Proteômica:} abundância de proteínas
    \item \textbf{Metabolômica:} concentrações de metabólitos
    \item \textbf{Fluxômica:} fluxos metabólicos
\end{itemize}

\column{0.5\textwidth}
\textbf{Desafios:}
\begin{itemize}
    \item Escalas diferentes
    \item Heterogeidade de dados
    \item Ruído experimental
    \item Dados faltantes
    \item Causalidade vs correlação
    \item Complexidade computacional
\end{itemize}
\end{columns}

\vspace{0.3cm}

\importante{
A integração efetiva requer métodos estatísticos robustos e modelagem matemática rigorosa.
}
\end{frame}

\begin{frame}{Caso de Estudo: Ciclo da Trealose em Levedura}
\begin{block}{Por que Trealose?}
\textit{Saccharomyces cerevisiae} utiliza trealose como:
\begin{itemize}
    \item Proteção contra estresse (calor, frio, dessecação)
    \item Reserva de carbono
    \item Sinalização celular
    \item Regulação do metabolismo
\end{itemize}
\end{block}

\vspace{0.3cm}

\begin{exampleblock}{Vantagens para Estudo Integrado}
\begin{itemize}
    \item Sistema bem caracterizado
    \item Via regulada em múltiplos níveis
    \item Dados multi-ômicos disponíveis
    \item Respostas mensuráveis
    \item Modelo para outros sistemas de estresse
\end{itemize}
\end{exampleblock}
\end{frame}

% ============================================================================
% SEÇÃO 2: INTEGRAÇÃO MULTI-ÔMICA
% ============================================================================
\section{Integração Multi-Ômica}

\begin{frame}{Visão Geral das Camadas Ômicas}
\begin{center}
\begin{tikzpicture}[scale=0.85]
    % Central dogma with omics layers
    \node[draw, rectangle, fill=roxodna!30, minimum width=2.5cm, minimum height=1cm] at (0,4) {Genômica};
    \node[draw, rectangle, fill=laranjacel!30, minimum width=2.5cm, minimum height=1cm] at (0,2.5) {Transcriptômica};
    \node[draw, rectangle, fill=verdeacido!30, minimum width=2.5cm, minimum height=1cm] at (0,1) {Proteômica};
    \node[draw, rectangle, fill=azulclaro!30, minimum width=2.5cm, minimum height=1cm] at (0,-0.5) {Metabolômica};

    % Arrows showing flow
    \draw[->, ultra thick] (0,3.5) -- (0,3.0) node[right, midway, font=\tiny] {transcrição};
    \draw[->, ultra thick] (0,2.0) -- (0,1.5) node[right, midway, font=\tiny] {tradução};
    \draw[->, ultra thick] (0,0.5) -- (0,0) node[right, midway, font=\tiny] {catálise};

    % Characteristics on right
    \node[right, font=\tiny, text width=3.5cm] at (2,4) {DNA: estático, variações\\SNPs, CNVs, mutações};
    \node[right, font=\tiny, text width=3.5cm] at (2,2.5) {mRNA: dinâmico, 10$^3$-10$^7$ cópias\\responde rapidamente};
    \node[right, font=\tiny, text width=3.5cm] at (2,1) {Proteínas: funcional, modificações\\10$^4$-10$^6$ moléculas};
    \node[right, font=\tiny, text width=3.5cm] at (2,-0.5) {Metabólitos: fenotípico\\10$^{-9}$-10$^{-3}$ M};
\end{tikzpicture}
\end{center}

\importante{
Cada camada adiciona informação complementar sobre o estado do sistema.
}
\end{frame}

\begin{frame}{Tipos e Características de Dados}
\begin{center}
\begin{tabular}{lllll}
\toprule
\textbf{Ômica} & \textbf{Técnica} & \textbf{Elementos} & \textbf{Faixa Dinâmica} & \textbf{Ruído} \\
\midrule
Genômica & WGS, SNP & $\sim$10$^4$ genes & Baixa & Baixo \\
Transcriptômica & RNA-seq & $\sim$10$^4$ mRNAs & Alta (10$^6$) & Médio \\
Proteômica & MS, Western & $\sim$10$^3$ prot. & Média (10$^4$) & Alto \\
Metabolômica & MS, NMR & $\sim$10$^2$ met. & Muito alta (10$^9$) & Alto \\
Fluxômica & $^{13}$C-MFA & $\sim$10$^2$ fluxos & Média & Médio \\
\bottomrule
\end{tabular}
\end{center}

\vspace{0.3cm}

\begin{exampleblock}{Implicações para Integração}
\begin{itemize}
    \item Normalização entre escalas
    \item Tratamento de valores faltantes
    \item Métodos estatísticos robustos
    \item Validação experimental
\end{itemize}
\end{exampleblock}
\end{frame}

\begin{frame}{Estratégias de Integração: Early, Intermediate, Late Fusion}
\begin{center}
\begin{tikzpicture}[scale=0.7, every node/.style={font=\tiny}]
    % Early fusion
    \node[font=\small, above] at (2,5) {\textbf{Early Fusion}};
    \node[draw, fill=roxodna!30] (g1) at (0,4) {Genômica};
    \node[draw, fill=laranjacel!30] (t1) at (1.5,4) {Transcr.};
    \node[draw, fill=verdeacido!30] (p1) at (3,4) {Proteômica};
    \node[draw, fill=azulclaro!30] (m1) at (4.5,4) {Metabol.};
    \node[draw, fill=yellow!30, text width=1.5cm, align=center] (int1) at (2.25,2.5) {Integração\\Completa};
    \node[draw, fill=gray!30, text width=1.5cm, align=center] (mod1) at (2.25,1) {Modelo};

    \draw[->, thick] (g1) -- (int1);
    \draw[->, thick] (t1) -- (int1);
    \draw[->, thick] (p1) -- (int1);
    \draw[->, thick] (m1) -- (int1);
    \draw[->, thick] (int1) -- (mod1);

    % Intermediate fusion
    \node[font=\small, above] at (8,5) {\textbf{Intermediate Fusion}};
    \node[draw, fill=roxodna!30] (g2) at (6.5,4) {Genômica};
    \node[draw, fill=laranjacel!30] (t2) at (8,4) {Transcr.};
    \node[draw, fill=verdeacido!30] (p2) at (9.5,4) {Proteômica};
    \node[draw, fill=yellow!30] (int2a) at (7.25,2.8) {Módulo 1};
    \node[draw, fill=yellow!30] (int2b) at (9.25,2.8) {Módulo 2};
    \node[draw, fill=gray!30, text width=1.5cm, align=center] (mod2) at (8.25,1) {Modelo};

    \draw[->, thick] (g2) -- (int2a);
    \draw[->, thick] (t2) -- (int2a);
    \draw[->, thick] (p2) -- (int2b);
    \draw[->, thick] (int2a) -- (mod2);
    \draw[->, thick] (int2b) -- (mod2);

    % Late fusion
    \node[font=\small, above] at (13.5,5) {\textbf{Late Fusion}};
    \node[draw, fill=roxodna!30, text width=1.2cm, align=center] (g3) at (12,4) {Genômica};
    \node[draw, fill=laranjacel!30, text width=1.2cm, align=center] (t3) at (13.5,4) {Transcr.};
    \node[draw, fill=verdeacido!30, text width=1.2cm, align=center] (p3) at (15,4) {Proteômica};
    \node[draw, fill=yellow!30] (m3a) at (12,2.8) {Modelo 1};
    \node[draw, fill=yellow!30] (m3b) at (13.5,2.8) {Modelo 2};
    \node[draw, fill=yellow!30] (m3c) at (15,2.8) {Modelo 3};
    \node[draw, fill=gray!30, text width=1.5cm, align=center] (final) at (13.5,1) {Consenso};

    \draw[->, thick] (g3) -- (m3a);
    \draw[->, thick] (t3) -- (m3b);
    \draw[->, thick] (p3) -- (m3c);
    \draw[->, thick] (m3a) -- (final);
    \draw[->, thick] (m3b) -- (final);
    \draw[->, thick] (m3c) -- (final);
\end{tikzpicture}
\end{center}
\end{frame}

\begin{frame}{Métodos Baseados em Correlação}
\begin{block}{Análise de Correlação Multi-Ômica}
\begin{equation*}
\rho_{XY} = \frac{\text{Cov}(X,Y)}{\sigma_X \sigma_Y} = \frac{E[(X-\mu_X)(Y-\mu_Y)]}{\sigma_X \sigma_Y}
\end{equation*}
\end{block}

\vspace{0.2cm}

\begin{columns}
\column{0.5\textwidth}
\textbf{Métodos:}
\begin{itemize}
    \item Correlação de Pearson (linear)
    \item Correlação de Spearman (monotônica)
    \item Distância de correlação
    \item Informação mútua
    \item Correlação parcial
    \item Correlação canônica (CCA)
\end{itemize}

\column{0.5\textwidth}
\textbf{Aplicações:}
\begin{itemize}
    \item mRNA-proteína
    \item Proteína-metabólito
    \item Expressão-fluxo
    \item Módulos co-regulados
\end{itemize}

\vspace{0.2cm}

\begin{alertblock}{Cuidado!}
Correlação $\neq$ Causalidade
\end{alertblock}
\end{columns}
\end{frame}

\begin{frame}{Integração Baseada em Redes}
\begin{columns}
\column{0.5\textwidth}
\begin{block}{Construção de Redes}
\begin{enumerate}
    \item \textbf{Nós:} biomoléculas (genes, proteínas, metabólitos)
    \item \textbf{Arestas:} interações, correlações, regulações
    \item \textbf{Pesos:} força da interação
\end{enumerate}
\end{block}

\textbf{Tipos de Redes:}
\begin{itemize}
    \item Co-expressão
    \item Interação proteína-proteína
    \item Regulação gênica
    \item Metabólica
    \item Multi-camada (multiplex)
\end{itemize}

\column{0.5\textwidth}
\begin{center}
\begin{tikzpicture}[scale=0.7]
    % Multi-layer network
    % Transcriptomics layer
    \node[draw, circle, fill=laranjacel!50, font=\tiny] (g1) at (0,3) {mRNA1};
    \node[draw, circle, fill=laranjacel!50, font=\tiny] (g2) at (1.5,3) {mRNA2};
    \node[draw, circle, fill=laranjacel!50, font=\tiny] (g3) at (3,3) {mRNA3};

    \draw[thick] (g1) -- (g2);
    \draw[thick] (g2) -- (g3);

    \node[left, font=\tiny] at (-0.5,3) {Transcriptômica};

    % Proteomics layer
    \node[draw, circle, fill=verdeacido!50, font=\tiny] (p1) at (0,1.5) {Prot1};
    \node[draw, circle, fill=verdeacido!50, font=\tiny] (p2) at (1.5,1.5) {Prot2};
    \node[draw, circle, fill=verdeacido!50, font=\tiny] (p3) at (3,1.5) {Prot3};

    \draw[thick] (p1) -- (p2);
    \draw[thick] (p2) -- (p3);
    \draw[thick] (p1) -- (p3);

    \node[left, font=\tiny] at (-0.5,1.5) {Proteômica};

    % Metabolomics layer
    \node[draw, circle, fill=azulclaro!50, font=\tiny] (m1) at (0,0) {Met1};
    \node[draw, circle, fill=azulclaro!50, font=\tiny] (m2) at (1.5,0) {Met2};
    \node[draw, circle, fill=azulclaro!50, font=\tiny] (m3) at (3,0) {Met3};

    \draw[thick] (m1) -- (m2);
    \draw[thick] (m2) -- (m3);

    \node[left, font=\tiny] at (-0.5,0) {Metabolômica};

    % Inter-layer connections
    \draw[dashed, azulescuro] (g1) -- (p1);
    \draw[dashed, azulescuro] (g2) -- (p2);
    \draw[dashed, azulescuro] (g3) -- (p3);
    \draw[dashed, azulescuro] (p1) -- (m1);
    \draw[dashed, azulescuro] (p2) -- (m2);
    \draw[dashed, azulescuro] (p3) -- (m3);
\end{tikzpicture}
\end{center}

\textbf{Análise:}
\begin{itemize}
    \item Centralidade
    \item Modularidade
    \item Comunidades
\end{itemize}
\end{columns}
\end{frame}

\begin{frame}{Abordagens de Machine Learning}
\begin{block}{Métodos Supervisionados}
\begin{itemize}
    \item \textbf{Random Forest:} classificação multi-ômica
    \item \textbf{SVM:} identificação de biomarcadores
    \item \textbf{Neural Networks:} predição de fenótipo
    \item \textbf{Elastic Net:} seleção de features
\end{itemize}
\end{block}

\begin{block}{Métodos Não-Supervisionados}
\begin{itemize}
    \item \textbf{PCA/PLS:} redução de dimensionalidade
    \item \textbf{t-SNE/UMAP:} visualização
    \item \textbf{Clustering:} identificação de subgrupos
    \item \textbf{NMF:} decomposição de matrizes
    \item \textbf{Autoencoders:} aprendizado de representações
\end{itemize}
\end{block}

\vspace{0.2cm}

\exemplo{MOFA (Multi-Omics Factor Analysis)}{
Decomposição fatorial para identificar fontes de variação compartilhadas e específicas entre camadas ômicas.
}
\end{frame}

\begin{frame}{Desafios: Heterogeneidade de Dados}
\begin{columns}
\column{0.5\textwidth}
\textbf{Tipos de Heterogeneidade:}
\begin{enumerate}
    \item \textbf{Escala:} valores absolutos vs relativos
    \item \textbf{Distribuição:} normal, log-normal, Poisson
    \item \textbf{Dimensionalidade:} 10$^2$ a 10$^5$ features
    \item \textbf{Esparsidade:} muitos zeros
    \item \textbf{Ruído:} técnico e biológico
\end{enumerate}

\column{0.5\textwidth}
\textbf{Soluções:}
\begin{itemize}
    \item Normalização apropriada
    \item Transformações (log, Box-Cox)
    \item Batch correction
    \item Imputação de valores faltantes
    \item Regularização
    \item Métodos robustos
\end{itemize}
\end{columns}

\vspace{0.3cm}

\begin{alertblock}{Normalização Multi-Ômica}
\begin{equation*}
X_{norm} = \frac{X - \mu}{\sigma} \quad \text{ou} \quad X_{norm} = \frac{X - \min(X)}{\max(X) - \min(X)}
\end{equation*}
\end{alertblock}
\end{frame}

\begin{frame}{Fluxo de Trabalho Multi-Ômico}
\begin{center}
\begin{tikzpicture}[node distance=1.2cm, scale=0.75, transform shape,
    box/.style={rectangle, draw, text width=3cm, text centered, rounded corners, minimum height=0.8cm, font=\small}]

    \node[box, fill=roxodna!30] (raw) {Dados Brutos\\Multi-Ômicos};
    \node[box, fill=laranjacel!30, below of=raw] (qc) {Controle de\\Qualidade};
    \node[box, fill=verdeacido!30, below of=qc] (norm) {Normalização\\e Filtro};
    \node[box, fill=azulclaro!30, below of=norm] (int) {Integração\\de Dados};
    \node[box, fill=yellow!30, below of=int] (anal) {Análise\\Estatística};
    \node[box, fill=gray!30, below of=anal] (model) {Modelagem\\de Sistemas};
    \node[box, fill=green!30, below of=model] (val) {Validação\\Experimental};

    \draw[->, thick] (raw) -- (qc);
    \draw[->, thick] (qc) -- (norm);
    \draw[->, thick] (norm) -- (int);
    \draw[->, thick] (int) -- (anal);
    \draw[->, thick] (anal) -- (model);
    \draw[->, thick] (model) -- (val);

    % Feedback loop
    \draw[->, thick, dashed, red] (val) -| ++(2,0) |- (model);

    % Annotations
    \node[right, font=\tiny, text width=3cm] at (4,-1.2) {Remover outliers\\FastQC, MultiQC};
    \node[right, font=\tiny, text width=3cm] at (4,-2.4) {Z-score, log\\batch correction};
    \node[right, font=\tiny, text width=3cm] at (4,-3.6) {Correlação\\Redes, ML};
    \node[right, font=\tiny, text width=3cm] at (4,-4.8) {Testes, FDR\\Enriquecimento};
    \node[right, font=\tiny, text width=3cm] at (4,-6.0) {ODEs, FBA\\Simulação};
    \node[right, font=\tiny, text width=3cm] at (4,-7.2) {qPCR, Western\\Ensaios enzimáticos};
\end{tikzpicture}
\end{center}
\end{frame}

\begin{frame}[fragile]{Exemplo Python: Integração de Dados}
\begin{lstlisting}[language=Python, caption=Integração multi-ômica com correlação, basicstyle=\ttfamily\scriptsize]
import pandas as pd
import numpy as np
from scipy import stats
import seaborn as sns
import matplotlib.pyplot as plt

# Carregar dados multi-omicos
transcriptomics = pd.read_csv('transcriptomics.csv', index_col=0)
proteomics = pd.read_csv('proteomics.csv', index_col=0)
metabolomics = pd.read_csv('metabolomics.csv', index_col=0)

# Normalizar cada dataset (z-score)
def normalize_zscore(df):
    return (df - df.mean()) / df.std()

trans_norm = normalize_zscore(transcriptomics)
prot_norm = normalize_zscore(proteomics)
metab_norm = normalize_zscore(metabolomics)

# Correlacao mRNA-proteina
correlation_matrix = pd.DataFrame(
    index=trans_norm.index,
    columns=prot_norm.index
)

for gene in trans_norm.index:
    for protein in prot_norm.index:
        r, p = stats.pearsonr(trans_norm.loc[gene],
                              prot_norm.loc[protein])
        if p < 0.05:  # significativo
            correlation_matrix.loc[gene, protein] = r

# Visualizar heatmap
plt.figure(figsize=(12, 10))
sns.heatmap(correlation_matrix.astype(float),
            cmap='RdBu_r', center=0, vmin=-1, vmax=1)
plt.title('Correlacao mRNA-Proteina')
\end{lstlisting}
\end{frame}

% ============================================================================
% SEÇÃO 3: CASO DE ESTUDO - CICLO DA TREALOSE
% ============================================================================
\section{Caso de Estudo: Ciclo da Trealose em Levedura}

\begin{frame}{Background Biológico: Resposta ao Estresse em \textit{S. cerevisiae}}
\begin{block}{Trealose como Protetor Celular}
\begin{itemize}
    \item Dissacarídeo não-redutor ($\alpha$-D-glicopiranosil-$\alpha$-D-glicopiranosídeo)
    \item Acumula durante estresse térmico, osmótico, oxidativo
    \item Estabiliza membranas e proteínas
    \item Pode atingir 20\% do peso seco celular
    \item Mobilizada quando condições melhoram
\end{itemize}
\end{block}

\vspace{0.3cm}

\begin{exampleblock}{Relevância Industrial}
\begin{itemize}
    \item Panificação: tolerância ao estresse
    \item Fermentação alcoólica: robustez
    \item Produção de biocombustíveis
    \item Biotecnologia: engenharia de robustez
\end{itemize}
\end{exampleblock}
\end{frame}

\begin{frame}{Via de Biossíntese e Degradação da Trealose}
\begin{center}
\begin{tikzpicture}[scale=0.85, every node/.style={font=\small}]
    % Synthesis pathway
    \node[draw, circle, fill=azulclaro!30] (g6p) at (0,3) {G6P};
    \node[draw, circle, fill=azulclaro!30] (udpg) at (0,1.5) {UDP-Glc};
    \node[draw, circle, fill=azulclaro!30] (t6p) at (3,1.5) {T6P};
    \node[draw, circle, fill=laranjacel!50, text width=1.5cm, align=center] (tre) at (6,1.5) {Trealose};

    % Enzymes
    \node[draw, rectangle, fill=verdeacido!50, font=\tiny] (tps1) at (1.5,2.2) {TPS1};
    \node[draw, rectangle, fill=verdeacido!50, font=\tiny] (tps2) at (4.5,2.2) {TPS2};
    \node[draw, rectangle, fill=red!50, font=\tiny] (nth1) at (6,0) {NTH1};

    % Reactions
    \draw[->, ultra thick, azulescuro] (g6p) -- (udpg);
    \draw[->, ultra thick, azulescuro] (udpg) -- (t6p) node[above, midway, font=\tiny] {+ G6P};
    \draw[->, ultra thick, azulescuro] (t6p) -- (tre);
    \draw[->, ultra thick, red] (tre) -- (6,0.3) node[right, midway, font=\tiny] {2 Glc};

    % Enzyme labels
    \draw[->, dashed] (tps1) -- (1.5,1.5);
    \draw[->, dashed] (tps2) -- (4.5,1.5);
    \draw[->, dashed] (nth1) -- (6,0.75);

    % Regulation
    \node[draw, ellipse, fill=yellow!30, font=\tiny] (stress) at (-2,1.5) {Estresse};
    \draw[->, thick, red] (stress) -- (udpg) node[above, midway, font=\tiny] {$\uparrow$ síntese};

    \node[draw, ellipse, fill=green!30, font=\tiny] (normal) at (8,1.5) {Recuperação};
    \draw[->, thick, blue] (normal) -- (tre) node[above, midway, font=\tiny] {$\uparrow$ degradação};

    % Annotations
    \node[below, font=\tiny] at (3,0.8) {T6P: trehalose-6-fosfato};
    \node[below, font=\tiny] at (3,0.3) {G6P: glicose-6-fosfato};
\end{tikzpicture}
\end{center}

\importante{
Via regulada transcricionalmente, pós-traducionalmente e alostericamente.
}
\end{frame}

\begin{frame}{Enzimas Chave: TPS1, TPS2, NTH1}
\begin{columns}
\column{0.5\textwidth}
\begin{block}{TPS1 (Trealose-6-P Sintase)}
\begin{itemize}
    \item Catalisa: UDP-Glc + G6P $\rightarrow$ T6P
    \item Subunidade catalítica
    \item Gene essencial
    \item Regulação: PKA, Snf1
\end{itemize}
\begin{equation*}
v_{TPS1} = \frac{V_{max}[UDP\text{-}Glc][G6P]}{K_m^{UDP} K_m^{G6P}}
\end{equation*}
\end{block}

\begin{block}{TPS2 (Trealose-6-P Fosfatase)}
\begin{itemize}
    \item Catalisa: T6P $\rightarrow$ Trealose + Pi
    \item Remove T6P (tóxico em excesso)
    \item Essencial para crescimento
\end{itemize}
\end{block}

\column{0.5\textwidth}
\begin{block}{NTH1 (Trealase Neutra)}
\begin{itemize}
    \item Catalisa: Trealose $\rightarrow$ 2 Glc
    \item Mobiliza reservas
    \item Ativa por PKA (cAMP)
    \item Inibida durante estresse
\end{itemize}
\begin{equation*}
v_{NTH1} = \frac{V_{max}[Tre]}{K_m + [Tre]} \cdot f(PKA)
\end{equation*}
\end{block}

\vspace{0.3cm}

\begin{alertblock}{Complexo TSC}
TPS1, TPS2 e subunidades regulatórias (TSL1, TPS3) formam complexo de 800 kDa
\end{alertblock}
\end{columns}
\end{frame}

\begin{frame}{Mecanismos Regulatórios}
\begin{center}
\begin{tikzpicture}[scale=0.8, every node/.style={font=\small}]
    % Transcriptional regulation
    \node[font=\footnotesize, above] at (2,4) {\textbf{Regulação Transcricional}};
    \node[draw, ellipse, fill=yellow!30] (stress) at (0,3) {Estresse};
    \node[draw, rectangle, fill=roxodna!30] (msn24) at (2,3) {Msn2/4};
    \node[draw, rectangle, fill=laranjacel!30] (tps1g) at (4,3) {\gene{TPS1}/\gene{TPS2}};

    \draw[->, thick] (stress) -- (msn24);
    \draw[->, thick] (msn24) -- (tps1g);

    % Post-translational regulation
    \node[font=\footnotesize, above] at (2,1.5) {\textbf{Regulação Pós-Traducional}};
    \node[draw, ellipse, fill=verdeacido!30] (glucose) at (0,0.5) {Glicose};
    \node[draw, rectangle, fill=azulclaro!30] (camp) at (2,0.5) {cAMP/PKA};
    \node[draw, rectangle, fill=yellow!30] (nth1p) at (4,0.5) {NTH1-P};

    \draw[->, thick] (glucose) -- (camp);
    \draw[->, thick] (camp) -- (nth1p) node[above, midway, font=\tiny] {fosforilação};

    % Allosteric regulation
    \node[font=\footnotesize, above] at (9,4) {\textbf{Regulação Alostérica}};
    \node[draw, circle, fill=azulclaro!30] (g6p2) at (7.5,3) {G6P};
    \node[draw, circle, fill=laranjacel!50] (t6p2) at (9,3) {T6P};
    \node[draw, rectangle, fill=verdeacido!30] (pfk) at (10.5,3) {PFK};

    \draw[->, thick, red] (t6p2) -- (pfk) node[above, midway, font=\tiny] {inibição};
    \draw[->, thick, blue] (g6p2) -- (pfk) node[above, midway, font=\tiny] {ativação};

    \node[below, font=\tiny, text width=3cm, align=center] at (9,2.2) {T6P inibe glicólise quando\\há acúmulo de trealose};
\end{tikzpicture}
\end{center}

\destaque{
Sistema de feedback multi-nível garante homeostase de trealose e G6P.
}
\end{frame}

\begin{frame}{Dados Experimentais: Multi-Ômicos}
\begin{block}{Dataset Integrado (Choque Térmico 37°C)}
\begin{center}
\begin{tabular}{llll}
\toprule
\textbf{Tempo (min)} & \textbf{mRNA TPS1} & \textbf{Proteína TPS1} & \textbf{[Trealose] mM} \\
\midrule
0 & 1.0 & 1.0 & 5 \\
15 & 3.2 & 1.2 & 8 \\
30 & 8.5 & 2.8 & 25 \\
60 & 12.0 & 5.5 & 85 \\
90 & 10.5 & 7.2 & 120 \\
120 & 6.8 & 6.8 & 130 \\
\bottomrule
\end{tabular}
\end{center}
\end{block}

\vspace{0.2cm}

\begin{exampleblock}{Observações}
\begin{itemize}
    \item mRNA responde rapidamente (15 min)
    \item Proteína segue com atraso (30-60 min)
    \item Trealose acumula gradualmente
    \item Saturação em 2h
\end{itemize}
\end{exampleblock}
\end{frame}

\begin{frame}{Construção do Modelo Matemático}
\begin{block}{Sistema de EDOs para Ciclo da Trealose}
\begin{align*}
\frac{d[mRNA_{TPS1}]}{dt} &= k_{trans} \cdot f(stress) - \delta_{mRNA}[mRNA_{TPS1}] \\
\frac{d[TPS1]}{dt} &= k_{transl}[mRNA_{TPS1}] - \delta_{prot}[TPS1] \\
\frac{d[T6P]}{dt} &= v_{TPS1} - v_{TPS2} \\
\frac{d[Tre]}{dt} &= v_{TPS2} - v_{NTH1}
\end{align*}

onde:
\begin{itemize}
    \item $f(stress) = \frac{stress^n}{K_d^n + stress^n}$ (função de Hill)
    \item $v_{TPS1} = \frac{V_{max}^{TPS1}[TPS1][G6P][UDP\text{-}Glc]}{(K_m^{G6P} + [G6P])(K_m^{UDP} + [UDP\text{-}Glc])}$
    \item $v_{NTH1} = \frac{V_{max}^{NTH1}[NTH1][Tre]}{K_m^{Tre} + [Tre]} \cdot f(PKA)$
\end{itemize}
\end{block}
\end{frame}

\begin{frame}{Estimação de Parâmetros}
\begin{columns}
\column{0.5\textwidth}
\textbf{Métodos de Estimação:}
\begin{enumerate}
    \item \textbf{Mínimos Quadrados:}
    \begin{equation*}
    \min_{\theta} \sum_i (y_i^{obs} - y_i^{model}(\theta))^2
    \end{equation*}

    \item \textbf{Maximum Likelihood:}
    \begin{equation*}
    \max_{\theta} \prod_i P(y_i^{obs}|\theta)
    \end{equation*}

    \item \textbf{Bayesiano (MCMC):}
    \begin{equation*}
    P(\theta|y) \propto P(y|\theta)P(\theta)
    \end{equation*}
\end{enumerate}

\column{0.5\textwidth}
\textbf{Restrições:}
\begin{itemize}
    \item $V_{max} > 0$
    \item $K_m > 0$
    \item $\delta > 0$
    \item Limites biológicos
\end{itemize}

\vspace{0.3cm}

\begin{exampleblock}{Ferramentas}
\begin{itemize}
    \item \texttt{scipy.optimize}
    \item \texttt{lmfit}
    \item \texttt{PyMC3}
    \item COPASI
    \item MATLAB simbiology
\end{itemize}
\end{exampleblock}
\end{columns}

\vspace{0.3cm}

\importante{
Usar dados multi-ômicos para restringir espaço de parâmetros.
}
\end{frame}

\begin{frame}{Validação do Modelo}
\begin{columns}
\column{0.5\textwidth}
\textbf{Critérios de Validação:}
\begin{enumerate}
    \item \textbf{Ajuste aos dados:} $R^2$, RMSE
    \item \textbf{Validação cruzada:} hold-out, k-fold
    \item \textbf{Predições independentes:} novas condições
    \item \textbf{Análise de sensibilidade:} parâmetros críticos
    \item \textbf{Testes experimentais:} mutantes, perturbações
\end{enumerate}

\column{0.5\textwidth}
\textbf{Métricas:}
\begin{itemize}
    \item $R^2 = 1 - \frac{SS_{res}}{SS_{tot}}$
    \item $RMSE = \sqrt{\frac{1}{n}\sum(y_i - \hat{y}_i)^2}$
    \item $AIC = 2k - 2\ln(L)$
    \item $BIC = k\ln(n) - 2\ln(L)$
\end{itemize}
\end{columns}

\vspace{0.3cm}

\begin{exampleblock}{Validação Experimental}
\begin{itemize}
    \item Testar mutantes: $\Delta$\gene{tps1}, $\Delta$\gene{nth1}
    \item Perturbações químicas: inibidores de PKA
    \item Diferentes condições de estresse
    \item Comparar predições com medições
\end{itemize}
\end{exampleblock}
\end{frame}

\begin{frame}{Predições e Verificação Experimental}
\begin{block}{Predições do Modelo}
\begin{enumerate}
    \item Aumento de 10x na expressão de TPS1 sob estresse leva a acúmulo 2x mais rápido de trealose
    \item Deleção de NTH1 resulta em acúmulo de trealose mesmo após estresse
    \item Inibição de PKA aumenta trealose basal em 3x
    \item Pré-adaptação ao estresse acelera resposta em 40\%
\end{enumerate}
\end{block}

\vspace{0.3cm}

\begin{exampleblock}{Verificação Experimental}
\begin{center}
\begin{tabular}{lcc}
\toprule
\textbf{Condição} & \textbf{Predição} & \textbf{Experimental} \\
\midrule
Superexpressão TPS1 & 2.1x rápido & 1.9x rápido \\
$\Delta$\gene{nth1} & 130 mM final & 125 mM final \\
Inibidor PKA & 3.2x basal & 3.0x basal \\
Pré-adaptação & 40\% $\downarrow$ tempo & 38\% $\downarrow$ tempo \\
\bottomrule
\end{tabular}
\end{center}
\end{exampleblock}

\destaque{
Concordância entre modelo e experimentos valida abordagem integrada!
}
\end{frame}

\begin{frame}[fragile]{Implementação Python: Modelo de Trealose}
\begin{lstlisting}[language=Python, caption=Simulação do ciclo da trealose, basicstyle=\ttfamily\scriptsize]
import numpy as np
from scipy.integrate import odeint
import matplotlib.pyplot as plt

def trehalose_model(y, t, params):
    """Modelo do ciclo da trealose"""
    mRNA_TPS1, TPS1, T6P, Tre = y

    k_trans, delta_mRNA, k_transl, delta_prot = params[:4]
    Vmax_TPS1, Km_G6P, Vmax_TPS2, Km_T6P = params[4:8]
    Vmax_NTH1, Km_Tre, stress_level = params[8:]

    # Constantes
    G6P = 5.0  # mM (assumido constante)
    UDP_Glc = 2.0  # mM

    # Regulacao transcricional (Hill function)
    n_hill = 4
    Kd_stress = 0.5
    f_stress = stress_level**n_hill / (Kd_stress**n_hill + stress_level**n_hill)

    # Taxas de reacao
    v_TPS1 = (Vmax_TPS1 * TPS1 * G6P * UDP_Glc) / ((Km_G6P + G6P) * (1.0 + UDP_Glc))
    v_TPS2 = (Vmax_TPS2 * T6P) / (Km_T6P + T6P)
    v_NTH1 = (Vmax_NTH1 * Tre) / (Km_Tre + Tre)

    # EDOs
    d_mRNA = k_trans * f_stress - delta_mRNA * mRNA_TPS1
    d_TPS1 = k_transl * mRNA_TPS1 - delta_prot * TPS1
    d_T6P = v_TPS1 - v_TPS2
    d_Tre = v_TPS2 - v_NTH1

    return [d_mRNA, d_TPS1, d_T6P, d_Tre]
\end{lstlisting}
\end{frame}

\begin{frame}[fragile]{Simulação e Ajuste aos Dados}
\begin{lstlisting}[language=Python, caption=Ajuste do modelo aos dados experimentais, basicstyle=\ttfamily\scriptsize]
# Parametros (estimados)
params = [
    5.0,   # k_trans
    0.1,   # delta_mRNA
    2.0,   # k_transl
    0.05,  # delta_prot
    10.0,  # Vmax_TPS1
    0.5,   # Km_G6P
    15.0,  # Vmax_TPS2
    0.2,   # Km_T6P
    8.0,   # Vmax_NTH1
    10.0,  # Km_Tre
    1.0    # stress_level (heat shock)
]

# Condicoes iniciais
y0 = [1.0, 1.0, 0.1, 5.0]  # mRNA, TPS1, T6P, Tre

# Tempo de simulacao
t = np.linspace(0, 120, 100)  # 0 a 120 min

# Resolver EDOs
sol = odeint(trehalose_model, y0, t, args=(params,))

# Plotar resultados
fig, axes = plt.subplots(2, 2, figsize=(12, 10))
axes[0,0].plot(t, sol[:,0], 'b-', label='mRNA TPS1')
axes[0,1].plot(t, sol[:,1], 'g-', label='Proteina TPS1')
axes[1,0].plot(t, sol[:,2], 'r-', label='T6P')
axes[1,1].plot(t, sol[:,3], 'orange', label='Trealose')

# Adicionar dados experimentais
t_exp = [0, 15, 30, 60, 90, 120]
tre_exp = [5, 8, 25, 85, 120, 130]
axes[1,1].plot(t_exp, tre_exp, 'ko', label='Experimental')
\end{lstlisting}
\end{frame}

% ============================================================================
% SEÇÃO 4: ANÁLISE DE VIAS E REDES INTEGRADAS
% ============================================================================
\section{Análise de Vias e Redes Integradas}

\begin{frame}{Análise de Enriquecimento de Vias (Pathway Enrichment)}
\definicao{Análise de Enriquecimento}{
Método estatístico para identificar se um conjunto de genes/proteínas está super-representado em vias biológicas ou processos específicos.
}

\begin{block}{Teste de Hipergeométrico}
\begin{equation*}
p = \frac{\binom{K}{k}\binom{N-K}{n-k}}{\binom{N}{n}}
\end{equation*}
onde:
\begin{itemize}
    \item $N$ = total de genes no genoma
    \item $K$ = genes na via de interesse
    \item $n$ = genes no conjunto (ex: diferencialmente expressos)
    \item $k$ = genes no conjunto que estão na via
\end{itemize}
\end{block}

\begin{exampleblock}{Ferramentas}
DAVID, Enrichr, g:Profiler, clusterProfiler, GSEA
\end{exampleblock}
\end{frame}

\begin{frame}{Gene Set Enrichment Analysis (GSEA)}
\begin{columns}
\column{0.5\textwidth}
\textbf{Vantagem sobre ORA:}
\begin{itemize}
    \item Usa ranking completo de genes
    \item Não requer threshold arbitrário
    \item Detecta mudanças coordenadas sutis
    \item Considera magnitude das mudanças
\end{itemize}

\vspace{0.3cm}

\textbf{Procedimento:}
\begin{enumerate}
    \item Rankear genes (fold-change, t-statistic)
    \item Calcular Enrichment Score (ES)
    \item Permutar para p-valor
    \item Corrigir para múltiplos testes (FDR)
\end{enumerate}

\column{0.5\textwidth}
\begin{center}
\begin{tikzpicture}[scale=0.6]
    % GSEA plot
    \draw[->] (0,0) -- (6,0) node[right] {Genes (ranked)};
    \draw[->] (0,0) -- (0,3) node[above] {ES};

    % Enrichment curve
    \draw[thick, azulescuro, domain=0:6, samples=50] plot (\x, {1.5 + 0.8*sin(180*\x/3)});

    % Gene set hits
    \foreach \x in {0.5, 1.0, 1.2, 1.5, 1.8, 2.2} {
        \draw[red, thick] (\x,0) -- (\x,0.3);
    }

    % Labels
    \node[below, font=\tiny] at (1,0) {Alta expr.};
    \node[below, font=\tiny] at (5,0) {Baixa expr.};
    \node[left, font=\tiny] at (0,2.3) {ES$_{max}$};

    % Enrichment peak
    \draw[dashed] (0,2.3) -- (1.8,2.3);
    \draw[dashed] (1.8,0) -- (1.8,2.3);
    \node[circle, fill=red, inner sep=1pt] at (1.8,2.3) {};
\end{tikzpicture}
\end{center}

\vspace{0.2cm}

\textbf{Interpretação:}
\begin{itemize}
    \item ES positivo: via ativada
    \item ES negativo: via reprimida
    \item Magnitude: força do enriquecimento
\end{itemize}
\end{columns}
\end{frame}

\begin{frame}{Reconstrução de Redes a partir de Multi-Ômicos}
\begin{block}{Abordagens de Reconstrução}
\begin{enumerate}
    \item \textbf{Correlação:} Pearson, Spearman
    \begin{itemize}
        \item Simples, rápido
        \item Não distingue correlação direta vs indireta
    \end{itemize}

    \item \textbf{Correlação Parcial:} controla co-variáveis
    \begin{equation*}
    \rho_{XY|Z} = \frac{\rho_{XY} - \rho_{XZ}\rho_{YZ}}{\sqrt{(1-\rho_{XZ}^2)(1-\rho_{YZ}^2)}}
    \end{equation*}

    \item \textbf{Informação Mútua:}
    \begin{equation*}
    I(X;Y) = \sum_x\sum_y p(x,y)\log\frac{p(x,y)}{p(x)p(y)}
    \end{equation*}

    \item \textbf{Graphical LASSO:} inferência de grafos Gaussianos
    \item \textbf{Redes Bayesianas:} aprendem causalidade
\end{enumerate}
\end{block}
\end{frame}

\begin{frame}{Detecção de Módulos e Clustering}
\begin{columns}
\column{0.5\textwidth}
\textbf{Métodos de Clustering:}
\begin{itemize}
    \item \textbf{K-means:} clusters esféricos
    \item \textbf{Hierárquico:} dendrograma
    \item \textbf{DBSCAN:} baseado em densidade
    \item \textbf{Louvain:} modularidade em grafos
    \item \textbf{Spectral:} autovalores da matriz de adjacência
\end{itemize}

\vspace{0.3cm}

\begin{block}{Modularidade}
\begin{equation*}
Q = \frac{1}{2m}\sum_{ij}\left[A_{ij} - \frac{k_i k_j}{2m}\right]\delta(c_i, c_j)
\end{equation*}
\end{block}

\column{0.5\textwidth}
\begin{center}
\begin{tikzpicture}[scale=0.6]
    % Network with modules
    % Module 1
    \node[draw, circle, fill=azulclaro!50, font=\tiny] (a1) at (0,2) {G1};
    \node[draw, circle, fill=azulclaro!50, font=\tiny] (a2) at (1,2.5) {G2};
    \node[draw, circle, fill=azulclaro!50, font=\tiny] (a3) at (1,1.5) {G3};

    \draw[thick] (a1) -- (a2);
    \draw[thick] (a1) -- (a3);
    \draw[thick] (a2) -- (a3);

    % Module 2
    \node[draw, circle, fill=verdeacido!50, font=\tiny] (b1) at (3,2) {G4};
    \node[draw, circle, fill=verdeacido!50, font=\tiny] (b2) at (4,2.5) {G5};
    \node[draw, circle, fill=verdeacido!50, font=\tiny] (b3) at (4,1.5) {G6};

    \draw[thick] (b1) -- (b2);
    \draw[thick] (b1) -- (b3);
    \draw[thick] (b2) -- (b3);

    % Module 3
    \node[draw, circle, fill=laranjacel!50, font=\tiny] (c1) at (2,0) {G7};
    \node[draw, circle, fill=laranjacel!50, font=\tiny] (c2) at (3,0) {G8};

    \draw[thick] (c1) -- (c2);

    % Inter-module edges
    \draw[dashed, thin] (a3) -- (b1);
    \draw[dashed, thin] (b3) -- (c1);

    % Labels
    \node[below, font=\tiny] at (0.5,3) {Módulo 1};
    \node[below, font=\tiny] at (3.5,3) {Módulo 2};
    \node[below, font=\tiny] at (2.5,-0.5) {Módulo 3};
\end{tikzpicture}
\end{center}

\textbf{Aplicações:}
\begin{itemize}
    \item Identificar processos biológicos
    \item Co-regulação
    \item Predição de função
\end{itemize}
\end{columns}

\vspace{0.3cm}

\exemplo{Módulo de Resposta ao Estresse}{
Genes co-expressos: HSP, TPS, CTT, SSA - todos regulados por Msn2/4
}
\end{frame}

\begin{frame}{Identificação de Reguladores e Hubs}
\begin{block}{Métricas de Centralidade}
\begin{itemize}
    \item \textbf{Grau (Degree):} número de conexões
    \begin{equation*}
    k_i = \sum_j A_{ij}
    \end{equation*}

    \item \textbf{Betweenness:} fração de caminhos mais curtos que passam pelo nó
    \begin{equation*}
    C_B(i) = \sum_{s \neq i \neq t} \frac{\sigma_{st}(i)}{\sigma_{st}}
    \end{equation*}

    \item \textbf{Closeness:} proximidade média a outros nós
    \begin{equation*}
    C_C(i) = \frac{n-1}{\sum_j d(i,j)}
    \end{equation*}

    \item \textbf{Eigenvector:} conectado a nós importantes
    \item \textbf{PageRank:} variante de eigenvector centralidade
\end{itemize}
\end{block}
\end{frame}

\begin{frame}{Análise Dinâmica de Redes}
\begin{columns}
\column{0.5\textwidth}
\textbf{Redes que Mudam no Tempo:}
\begin{itemize}
    \item Topologia varia com condições
    \item Arestas aparecem/desaparecem
    \item Pesos mudam
    \item Módulos se reorganizam
\end{itemize}

\vspace{0.3cm}

\textbf{Métodos:}
\begin{itemize}
    \item Time-lagged correlation
    \item Dynamic Bayesian Networks
    \item Graphical Granger causality
    \item Sliding window analysis
\end{itemize}

\column{0.5\textwidth}
\begin{center}
\begin{tikzpicture}[scale=0.6]
    % Time series of networks
    \node[font=\tiny] at (0,3) {t = 0};
    \node[draw, circle, fill=azulclaro!50, minimum size=0.3cm] (t0a) at (0,2) {};
    \node[draw, circle, fill=azulclaro!50, minimum size=0.3cm] (t0b) at (0.8,2) {};
    \node[draw, circle, fill=azulclaro!50, minimum size=0.3cm] (t0c) at (0.4,1.2) {};
    \draw[thick] (t0a) -- (t0b);

    \node[font=\tiny] at (2.5,3) {t = 30 min};
    \node[draw, circle, fill=laranjacel!50, minimum size=0.4cm] (t1a) at (2.5,2) {};
    \node[draw, circle, fill=laranjacel!50, minimum size=0.3cm] (t1b) at (3.3,2) {};
    \node[draw, circle, fill=laranjacel!50, minimum size=0.3cm] (t1c) at (2.9,1.2) {};
    \draw[thick] (t1a) -- (t1b);
    \draw[thick] (t1a) -- (t1c);
    \draw[thick] (t1b) -- (t1c);

    \node[font=\tiny] at (5,3) {t = 60 min};
    \node[draw, circle, fill=verdeacido!50, minimum size=0.6cm] (t2a) at (5,2) {};
    \node[draw, circle, fill=verdeacido!50, minimum size=0.5cm] (t2b) at (5.8,2) {};
    \node[draw, circle, fill=verdeacido!50, minimum size=0.4cm] (t2c) at (5.4,1.2) {};
    \draw[ultra thick] (t2a) -- (t2b);
    \draw[ultra thick] (t2a) -- (t2c);
    \draw[thick] (t2b) -- (t2c);

    % Arrows showing time
    \draw[->, thick] (1,1.5) -- (1.8,1.5);
    \draw[->, thick] (3.5,1.5) -- (4.3,1.5);

    \node[below, font=\tiny, text width=2cm, align=center] at (3,-0.2) {Resposta ao estresse térmico};
\end{tikzpicture}
\end{center}

\vspace{0.2cm}

\importante{
Redes dinâmicas revelam sequência temporal de eventos regulatórios.
}
\end{columns}
\end{frame}

\begin{frame}[fragile]{Exemplo Python: Análise de Redes com NetworkX}
\begin{lstlisting}[language=Python, caption=Construção e análise de rede multi-ômica, basicstyle=\ttfamily\scriptsize]
import networkx as nx
import pandas as pd
import matplotlib.pyplot as plt
from scipy.stats import pearsonr

# Construir rede de correlacao
def build_correlation_network(data, threshold=0.7):
    """Cria rede baseada em correlacoes"""
    G = nx.Graph()

    # Adicionar nos
    for gene in data.index:
        G.add_node(gene)

    # Adicionar arestas (correlacoes significativas)
    for i, gene1 in enumerate(data.index):
        for gene2 in data.index[i+1:]:
            r, p = pearsonr(data.loc[gene1], data.loc[gene2])
            if abs(r) > threshold and p < 0.05:
                G.add_edge(gene1, gene2, weight=abs(r))

    return G

# Analisar centralidade
def analyze_network(G):
    """Calcula metricas de centralidade"""
    degree_cent = nx.degree_centrality(G)
    between_cent = nx.betweenness_centrality(G)
    close_cent = nx.closeness_centrality(G)

    # Identificar hubs (top 10% degree)
    threshold = sorted(degree_cent.values())[-len(G)//10]
    hubs = [node for node, cent in degree_cent.items()
            if cent >= threshold]

    return degree_cent, between_cent, close_cent, hubs
\end{lstlisting}
\end{frame}

\begin{frame}[fragile]{Análise de Enriquecimento com Python}
\begin{lstlisting}[language=Python, caption=GSEA com gseapy, basicstyle=\ttfamily\scriptsize]
import gseapy as gp
import pandas as pd

# Dados de expressao diferencial
de_genes = pd.read_csv('diff_expr.csv')
de_genes = de_genes.sort_values('log2FC', ascending=False)

# Criar ranking
gene_rank = de_genes.set_index('gene')['log2FC'].to_dict()

# Executar GSEA pre-ranked
gsea_results = gp.prerank(
    rnk=gene_rank,
    gene_sets='KEGG_2021_Human',  # banco de vias
    processes=4,
    permutation_num=1000,
    outdir='gsea_output',
    format='png',
    seed=42
)

# Vias significativamente enriquecidas
enriched_pathways = gsea_results.res2d[
    gsea_results.res2d['FDR q-val'] < 0.05
]

print(f"Vias enriquecidas: {len(enriched_pathways)}")
print(enriched_pathways[['Term', 'ES', 'NES', 'FDR q-val']].head(10))

# Over-representation analysis (ORA)
enrich_ora = gp.enrichr(
    gene_list=de_genes[de_genes['padj'] < 0.05]['gene'].tolist(),
    gene_sets='GO_Biological_Process_2021',
    organism='Yeast',
    outdir='enrichr_output'
)
\end{lstlisting}
\end{frame}

% ============================================================================
% SEÇÃO 5: MODELAGEM DE SISTEMAS INTEGRADOS
% ============================================================================
\section{Modelagem de Sistemas Integrados}

\begin{frame}{Modelos de Escala Genômica com Restrições Ômicas}
\definicao{Genome-Scale Model (GEM)}{
Reconstrução matemática do metabolismo completo de um organismo, incluindo todas as reações conhecidas, genes e estequiometria.
}

\begin{block}{Integração de Dados de Expressão}
\begin{itemize}
    \item \textbf{E-Flux:} usar níveis de expressão como proxy para capacidades enzimáticas
    \begin{equation*}
    v_{max,i} \propto \text{expression}_i
    \end{equation*}

    \item \textbf{GIMME:} minimizar diferenças entre fluxos e expressão
    \begin{equation*}
    \min \sum_i |v_i - v_i^{target}| \quad \text{s.t.} \quad Sv = 0
    \end{equation*}

    \item \textbf{iMAT:} maximizar consistência com dados binários (on/off)
\end{itemize}
\end{block}
\end{frame}

\begin{frame}{Constraint-Based Modeling com Transcriptômica}
\begin{columns}
\column{0.5\textwidth}
\textbf{Análise de Balanço de Fluxos (FBA):}
\begin{align*}
\max_v \quad & c^T v \\
\text{s.t.} \quad & Sv = 0 \\
& lb \leq v \leq ub
\end{align*}

\vspace{0.2cm}

\textbf{Com dados de expressão:}
\begin{itemize}
    \item Ajustar $lb$ e $ub$ baseado em níveis de mRNA
    \item Genes altamente expressos: $ub$ alto
    \item Genes não-expressos: $v = 0$
    \item Genes intermediários: $ub$ proporcional
\end{itemize}

\column{0.5\textwidth}
\begin{center}
\begin{tikzpicture}[scale=0.6]
    % FBA concept
    \node[font=\tiny, above] at (2,4) {\textbf{Espaço de Soluções}};

    % Feasible region
    \fill[azulclaro!20] (0,0) -- (4,0) -- (4,3) -- (0,3) -- cycle;

    % Expression constraints
    \fill[verdeacido!40] (0.5,0.5) -- (3.5,0.5) -- (3.5,2.5) -- (0.5,2.5) -- cycle;

    % Optimal solution
    \node[circle, fill=red, inner sep=2pt, label=right:{Ótimo}] at (3,2) {};

    % Axes
    \draw[->] (0,0) -- (4.5,0) node[right, font=\tiny] {$v_1$};
    \draw[->] (0,0) -- (0,3.5) node[above, font=\tiny] {$v_2$};

    % Labels
    \node[below, font=\tiny, text width=2cm, align=center] at (2,-0.5) {FBA padrão};
    \node[below, font=\tiny, text width=2cm, align=center] at (2,1.5) {Com expressão};

    % Arrows
    \draw[->, thick, red] (3.8,2.5) -- (3,2);
\end{tikzpicture}
\end{center}

\vspace{0.2cm}

\importante{
Dados de expressão restringem espaço de soluções, melhorando predições.
}
\end{columns}
\end{frame}

\begin{frame}{Modelos Cinéticos com Proteômica/Metabolômica}
\begin{block}{Integração Multi-Nível}
\begin{itemize}
    \item \textbf{Proteômica:} concentrações enzimáticas ($[E]$)
    \item \textbf{Metabolômica:} concentrações de substratos/produtos ($[S]$, $[P]$)
    \item \textbf{Fluxômica:} fluxos metabólicos ($v$)
\end{itemize}
\end{block}

\vspace{0.2cm}

\begin{block}{Modelo Cinético Integrado}
Para cada reação $i$:
\begin{equation*}
v_i = [E_i] \cdot k_{cat,i} \cdot \frac{[S_i]}{K_m^S + [S_i]} \cdot \left(1 - \frac{[P_i]}{[S_i]K_{eq,i}}\right)
\end{equation*}

Sistema de EDOs:
\begin{equation*}
\frac{d[S_i]}{dt} = \sum_j s_{ij}v_j
\end{equation*}

onde $s_{ij}$ é o coeficiente estequiométrico.
\end{block}

\vspace{0.2cm}

\destaque{
Dados proteômicos e metabolômicos fornecem valores absolutos para parametrização!
}
\end{frame}

\begin{frame}{Abordagens de Modelagem Híbrida}
\begin{columns}
\column{0.5\textwidth}
\textbf{Combinando Diferentes Formalismos:}
\begin{enumerate}
    \item \textbf{FBA + ODEs:}
    \begin{itemize}
        \item FBA para metabolismo central
        \item ODEs para regulação
    \end{itemize}

    \item \textbf{Estocástico + Determinístico:}
    \begin{itemize}
        \item Gillespie para expressão gênica
        \item ODEs para metabolismo
    \end{itemize}

    \item \textbf{Multi-escala:}
    \begin{itemize}
        \item Molecular (ms)
        \item Celular (min-h)
        \item População (h-dias)
    \end{itemize}
\end{enumerate}

\column{0.5\textwidth}
\begin{center}
\begin{tikzpicture}[scale=0.7]
    % Hybrid model structure
    \node[draw, rectangle, fill=roxodna!30, text width=2cm, align=center, font=\tiny] at (0,3) {Regulação\\Gênica\\(ODE)};

    \node[draw, rectangle, fill=laranjacel!30, text width=2cm, align=center, font=\tiny] at (0,1.5) {Tradução\\(Estocástico)};

    \node[draw, rectangle, fill=verdeacido!30, text width=2cm, align=center, font=\tiny] at (0,0) {Metabolismo\\(FBA)};

    % Connections
    \draw[<->, thick] (0,2.7) -- (0,1.8);
    \draw[<->, thick] (0,1.2) -- (0,0.3);

    % Time scales
    \node[right, font=\tiny] at (2.5,3) {Escala: min-h};
    \node[right, font=\tiny] at (2.5,1.5) {Escala: seg-min};
    \node[right, font=\tiny] at (2.5,0) {Escala: steady-state};
\end{tikzpicture}
\end{center}

\vspace{0.2cm}

\begin{exampleblock}{Vantagens}
\begin{itemize}
    \item Captura múltiplas escalas
    \item Usa melhor formalismo para cada nível
    \item Mais realista
\end{itemize}
\end{exampleblock}
\end{columns}
\end{frame}

\begin{frame}{Validação com Dados Independentes}
\begin{block}{Estratégias de Validação}
\begin{enumerate}
    \item \textbf{Hold-out:} treinar em 70\%, testar em 30\%
    \item \textbf{Validação cruzada:} k-fold (k=5 ou 10)
    \item \textbf{Leave-one-out:} para datasets pequenos
    \item \textbf{Validação temporal:} treinar em tempo $t$, validar em $t+\Delta t$
    \item \textbf{Validação experimental:} perturbações independentes
\end{enumerate}
\end{block}

\vspace{0.3cm}

\begin{columns}
\column{0.5\textwidth}
\textbf{Tipos de Dados para Validação:}
\begin{itemize}
    \item Mutantes deletados
    \item Superexpressões
    \item Diferentes condições nutricionais
    \item Resposta a drogas
    \item Dados de literatura
\end{itemize}

\column{0.5\textwidth}
\textbf{Métricas de Performance:}
\begin{itemize}
    \item Acurácia de predição
    \item Correlação observado vs predito
    \item Erro médio absoluto
    \item Sensibilidade/especificidade
    \item AUC-ROC
\end{itemize}
\end{columns}
\end{frame}

\begin{frame}[fragile]{Exemplo Python: FBA Integrado com Transcriptômica}
\begin{lstlisting}[language=Python, caption=FBA com restrições de expressão usando COBRApy, basicstyle=\ttfamily\scriptsize]
import cobra
import pandas as pd

# Carregar modelo genome-scale
model = cobra.io.read_sbml_model('iMM904.xml')  # S. cerevisiae

# Carregar dados de expressao (RNA-seq)
expression_data = pd.read_csv('expression_stress.csv', index_col=0)
expression_data = expression_data / expression_data.max()  # normalizar 0-1

# Integrar dados de expressao (metodo E-Flux)
def integrate_expression(model, expression, percentile=25):
    """Ajusta bounds baseado em expressao"""
    for reaction in model.reactions:
        # Encontrar genes associados
        genes = reaction.genes

        if len(genes) > 0:
            # Media de expressao dos genes
            gene_expr = [expression.get(g.id, 0.5) for g in genes]
            avg_expr = sum(gene_expr) / len(gene_expr)

            # Ajustar upper bound
            if avg_expr < percentile/100:
                # Gene pouco expresso: bloquear reacao
                reaction.upper_bound = 0
                reaction.lower_bound = 0
            else:
                # Escalar bound pela expressao
                original_ub = reaction.upper_bound
                if original_ub > 0:
                    reaction.upper_bound = original_ub * avg_expr

# Aplicar restricoes de expressao
integrate_expression(model, expression_data)

# Executar FBA
solution = model.optimize()
print(f"Taxa de crescimento predita: {solution.objective_value:.3f}")
\end{lstlisting}
\end{frame}

\begin{frame}[fragile]{Modelo Híbrido: Regulação + Metabolismo}
\begin{lstlisting}[language=Python, caption=Modelo híbrido simples, basicstyle=\ttfamily\scriptsize]
import numpy as np
from scipy.integrate import odeint
import cobra

def hybrid_model(y, t, model, params):
    """Modelo hibrido: ODE para regulacao + FBA para metabolismo"""
    mRNA, protein = y
    k_trans, k_transl, delta_mRNA, delta_prot = params

    # Parte 1: Regulacao genica (ODE)
    stress = 1.0 if t < 60 else 0.0  # estresse nos primeiros 60 min

    dmRNA_dt = k_trans * stress - delta_mRNA * mRNA
    dprotein_dt = k_transl * mRNA - delta_prot * protein

    # Parte 2: Ajustar metabolismo baseado em nivel proteico
    # Exemplo: protein = TPS1, afeta fluxo de trealose
    for rxn in model.reactions:
        if 'TPS1' in rxn.id:
            # Escalar capacidade pela abundancia proteica
            rxn.upper_bound = 10.0 * protein

    # Resolver FBA
    try:
        solution = model.optimize()
        growth_rate = solution.objective_value
        trehalose_flux = solution.fluxes.get('TPS1_rxn', 0.0)
    except:
        growth_rate = 0.0
        trehalose_flux = 0.0

    # Feedback: trealose afeta transcricao
    feedback = trehalose_flux / 10.0
    dmRNA_dt += feedback

    return [dmRNA_dt, dprotein_dt]

# Condicoes iniciais e parametros
y0 = [1.0, 1.0]
params = [5.0, 2.0, 0.1, 0.05]
t = np.linspace(0, 120, 100)

# Simular
# sol = odeint(hybrid_model, y0, t, args=(model, params))
\end{lstlisting}
\end{frame}

% ============================================================================
% SEÇÃO 6: EXERCÍCIOS
% ============================================================================
\section{Exercícios}

\begin{frame}{Exercícios - Parte 1: Correlação Multi-Ômica}
\begin{block}{Questão 1: Análise de Correlação}
Dados de transcriptômica e proteômica para 5 genes sob condição de estresse:
\begin{center}
\begin{tabular}{lcc}
\toprule
\textbf{Gene} & \textbf{log$_2$(mRNA)} & \textbf{log$_2$(Proteína)} \\
\midrule
\gene{TPS1} & 3.2 & 1.8 \\
\gene{HSP104} & 4.5 & 2.1 \\
\gene{CTT1} & 2.8 & 1.5 \\
\gene{SSA1} & 3.7 & 2.0 \\
\gene{GLK1} & 0.5 & 0.3 \\
\bottomrule
\end{tabular}
\end{center}
\begin{enumerate}
    \item Calcule a correlação de Pearson entre mRNA e proteína
    \item É estatisticamente significativa?
    \item Interprete biologicamente o resultado
\end{enumerate}
\end{block}
\end{frame}

\begin{frame}{Exercícios - Parte 2: Enriquecimento de Vias}
\begin{block}{Questão 2: Teste de Hipergeométrico}
Em um experimento de estresse oxidativo, 150 genes foram diferencialmente expressos de um total de 6000 genes no genoma. A via de "Resposta ao Estresse Oxidativo" contém 80 genes, e 25 deles estão no conjunto de genes diferencialmente expressos.
\begin{enumerate}
    \item Calcule o p-valor usando o teste hipergeométrico
    \item A via está significativamente enriquecida ($\alpha = 0.05$)?
    \item Se aplicarmos correção de Bonferroni para 300 vias testadas, ainda é significativa?
\end{enumerate}
\end{block}

\begin{block}{Questão 3: GSEA}
\begin{enumerate}
    \item Explique a diferença entre ORA e GSEA
    \item Por que GSEA pode detectar mudanças sutis que ORA perde?
    \item Quando você usaria cada método?
\end{enumerate}
\end{block}
\end{frame}

\begin{frame}{Exercícios - Parte 3: Construção de Modelos}
\begin{alertblock}{Questão 4: Projeto - Modelagem Integrada}
Construa um modelo integrado para a via de glicólise em \textit{S. cerevisiae}:
\begin{enumerate}
    \item Baixe dados de RNA-seq, proteômica e metabolômica do GEO
    \item Identifique genes/proteínas/metabólitos da glicólise
    \item Construa modelo cinético simplificado (3-5 reações chave)
    \item Estime parâmetros usando dados experimentais
    \item Valide o modelo com dados de mutantes
    \item Simule resposta a mudança de glicose extracelular
\end{enumerate}

\textbf{Entregáveis:}
\begin{itemize}
    \item Código Python documentado
    \item Gráficos de ajuste e validação
    \item Análise de sensibilidade de parâmetros
    \item Relatório interpretando resultados biologicamente
\end{itemize}
\end{alertblock}
\end{frame}

\begin{frame}{Exercícios - Parte 4: Análise de Redes}
\begin{block}{Questão 5: Construção de Rede de Co-expressão}
Usando dados de RNA-seq de séries temporais (0, 15, 30, 60, 120 min após estresse):
\begin{enumerate}
    \item Calcule matriz de correlação entre todos os pares de genes
    \item Construa rede mantendo correlações $|r| > 0.8$ e $p < 0.01$
    \item Identifique módulos usando algoritmo de Louvain
    \item Calcule centralidade de grau, betweenness e closeness
    \item Identifique top 10 hubs
    \item Faça análise de enriquecimento GO para cada módulo
\end{enumerate}
\end{block}

\begin{block}{Questão 6: FBA com Restrições}
\begin{enumerate}
    \item Carregue modelo genome-scale de \textit{S. cerevisiae} (iMM904)
    \item Integre dados de expressão usando método E-Flux
    \item Compare taxa de crescimento predita com/sem restrições
    \item Identifique reações essenciais sob estresse
\end{enumerate}
\end{block}
\end{frame}

% ============================================================================
% SEÇÃO 7: REFERÊNCIAS
% ============================================================================
\section{Referências}

\begin{frame}{Referências Principais}
\begin{thebibliography}{99}
\bibitem{klipp} Klipp E. et al. (2016). \textit{Systems Biology: A Textbook}, 2nd ed. Wiley-VCH.

\bibitem{palsson} Palsson B.O. (2015). \textit{Systems Biology: Constraint-based Reconstruction and Analysis}. Cambridge University Press.

\bibitem{ideker} Ideker T., Galitski T., Hood L. (2001). A new approach to decoding life: systems biology. \textit{Ann Rev Genomics Hum Genet}, 2:343-372.

\bibitem{hasin} Hasin Y. et al. (2017). Multi-omics approaches to disease. \textit{Genome Biology}, 18:83.

\bibitem{subramanian} Subramanian A. et al. (2005). Gene set enrichment analysis: a knowledge-based approach. \textit{PNAS}, 102(43):15545-15550.
\end{thebibliography}
\end{frame}

\begin{frame}{Leitura Complementar}
\begin{block}{Reviews Multi-Ômicos}
\begin{itemize}
    \item Ritchie M.D. et al. (2015). Methods of integrating data to uncover genotype-phenotype interactions. \textit{Nat Rev Genet}.
    \item Huang S. et al. (2017). More is better: recent progress in multi-omics data integration methods. \textit{Front Genet}.
    \item Cavill R. et al. (2016). Transcriptomic and metabolomic data integration. \textit{Brief Bioinform}.
\end{itemize}
\end{block}

\begin{block}{Bancos de Dados e Ferramentas}
\begin{itemize}
    \item \textbf{COBRApy:} \url{https://opencobra.github.io/cobrapy/}
    \item \textbf{NetworkX:} \url{https://networkx.org}
    \item \textbf{GSEApy:} \url{https://gseapy.readthedocs.io}
    \item \textbf{Enrichr:} \url{https://maayanlab.cloud/Enrichr/}
    \item \textbf{STRING:} \url{https://string-db.org}
    \item \textbf{KEGG:} \url{https://www.genome.jp/kegg/}
    \item \textbf{Reactome:} \url{https://reactome.org}
\end{itemize}
\end{block}
\end{frame}

\begin{frame}{Recursos Online}
\begin{block}{Tutoriais e Cursos}
\begin{itemize}
    \item \textbf{Systems Biology Course (MIT):} \url{https://ocw.mit.edu}
    \item \textbf{Network Analysis in Systems Biology:} \url{https://www.coursera.org}
    \item \textbf{Multi-omics Integration:} \url{https://www.ebi.ac.uk/training}
\end{itemize}
\end{block}

\begin{block}{Pacotes Python/R}
\begin{itemize}
    \item \textbf{Python:} cobra, networkx, gseapy, scanpy, scipy, scikit-learn
    \item \textbf{R:} clusterProfiler, WGCNA, mixOmics, limma, DESeq2
\end{itemize}
\end{block}

\begin{block}{Livros Especializados}
\begin{itemize}
    \item Alon U. (2019). \textit{An Introduction to Systems Biology}, 2nd ed.
    \item Voit E.O. (2017). \textit{A First Course in Systems Biology}, 2nd ed.
    \item Covert M.W. (2014). \textit{Fundamentals of Systems Biology}.
\end{itemize}
\end{block}
\end{frame}

% ============================================================================
% SLIDE FINAL
% ============================================================================
\begin{frame}
\begin{center}
{\Huge Obrigado!}

\vspace{1cm}

{\Large Capítulo 11: Integrated Analysis}

\vspace{1cm}

\textbf{Tópicos Abordados:}
\begin{itemize}
    \item Integração multi-ômica
    \item Caso de estudo: ciclo da trealose
    \item Análise de vias e redes
    \item Modelagem de sistemas integrados
\end{itemize}

\vspace{1cm}

\textit{Dúvidas? Perguntas?}
\end{center}
\end{frame}

\end{document}
